\subsection{Analisis Kinerja Sistem Elektronik}
\begin{itemize}
    \item Bahas efisiensi dan stabilitas \textit{power board}.
    \item Tinjau apakah integrasi STM32 dan Jetson Nano dalam satu PCB berhasil menurunkan ukuran dan mempercepat \textit{booting}.
    \item Diskusikan potensi \textit{noise}, isolasi daya, dan \textit{reliability} sistem.
\end{itemize}

\subsection{Analisis Kinerja Firmware dan Komunikasi}
\begin{itemize}
    \item Evaluasi kestabilan komunikasi servo (latensi, \textit{packet loss}).
    \item Analisis kecepatan respon antara Jetson dan STM32 (waktu \textit{round-trip} data).
    \item Bandingkan performa \textit{Sync Write/Read} dengan metode konvensional.
\end{itemize}

\subsection{Analisis Kinerja Mekanik dan Lengan 6-DOF}
\begin{itemize}
    \item Evaluasi jangkauan gerak dan ketepatan posisi berdasarkan hasil pengujian.
    \item Bandingkan dengan versi 3-DOF sebelumnya (berdasarkan batasan masalah Bab~1.3).
    \item Diskusikan kelebihan dan kendala dari desain baru.
\end{itemize}

\subsection{Analisis Kinerja Gerak dan Kinematika}
\begin{itemize}
    \item Analisis perbandingan hasil simulasi dan implementasi nyata.
    \item Evaluasi akurasi \textit{inverse kinematic} dan \textit{trajectory planning}.
    \item Pembahasan tentang kehalusan (\textit{smoothness}) dan stabilitas gerakan.
\end{itemize}

\subsection{Analisis Keseluruhan Sistem dan Ketercapaian Tujuan}
\begin{itemize}
    \item Bandingkan hasil yang diperoleh dengan tujuan penelitian (Bab~1.4).
    \item Diskusikan keberhasilan integrasi, efisiensi, dan peningkatan fungsi manipulasi.
    \item Uraikan limitasi sistem dan potensi pengembangan di masa depan.
\end{itemize}